\documentclass{beamer}

\usepackage{kotex}

\usepackage{fontspec}
\setmainfont{Pretendard Regular}

\title{Phoenix}
\author{2019320057 장창서}
\date{2025-10-01}

\begin{document}

\frame{\titlepage}

\begin{frame}
\frametitle{개괄}

\begin{itemize}
\item Elixir 배우기
\item Phoenix 배우기
\end{itemize}
\end{frame}

\begin{frame}
\frametitle{Elixir 배우기}

1주차: 함수형 기초
Getting Started: Introduction, Basic Types, Pattern Matching
과제: 리스트 길이 세기, 튜플 매칭 함수

2주차: 데이터와 함수
Getting Started: Modules, Functions, Control Structures
과제: Enum/Map 활용 파이프라인

3주차: 고급 함수 활용
Getting Started: Recursion, Mix and OTP (기초)
과제: map/filter 재귀 구현, 피보나치

4주차: 프로세스 기초
Getting Started: Processes
과제: 메시지 기반 카운터

5주차: OTP 심화
Getting Started: IO and the File System, Mix and OTP (심화)
과제: GenServer 기반 저장소, PubSub 브로드캐스트
\end{frame}

\begin{frame}
\frametitle{Phoenix 배우기}

6주차: Phoenix 구조와 LiveView 입문
Phoenix Guides: Up and Running, Routing, LiveView: Introduction
과제: 카운터 LiveView

7주차: 상태 관리와 LiveView 연동
Phoenix Guides: LiveView: Events, LiveView: Streams
과제: TodoStore + LiveView 연동

8주차: 실시간 동기화와 마무리
Phoenix Guides: PubSub, LiveView: Lifecycle
과제: 두 페이지 동기화 검증, 발표
\end{frame}

\begin{frame}
\frametitle{세션 형식}

\begin{itemize}
\item 기본적으로 제가 세션을 진행하고
\item 과제는 모두 GitHub PR로 제출
\item 과제하면서 어떤 부분이 어려웠는지 어떤 부분을 배웠는지 얘기나눠볼거 같습니다.
\end{itemize}
\end{frame}

\end{document}

